\documentclass[10.5pt,a4paper]{article}
\usepackage[T1]{fontenc}
\usepackage[margin=2cm]{geometry}
\usepackage{parskip}
\begin{document}
\pagestyle{empty}

\begin{flushright}
[Date]
\end{flushright}

The Editor-in-Chief\\
\textit{[Journal name]}

\medskip
\textbf{Subject: Submission of the manuscript ``Fractional-Memory Piezoelectric--Flexoelectric Strain-Gradient Strips with Rectangular-Window Second Moments: a $C^{1}$ Finite Element Framework''}

\medskip
Dear Editor,

We submit the above manuscript for consideration as an original research article in \textit{[Journal name]}.

The paper develops a $C^{1}$ Bogner--Fox--Schmit finite element formulation for a piezoelectric--flexoelectric strain-gradient strip in which the mechanical stiffness carries fractional (Caputo) memory. The gradient-energy coefficients of the adopted local model are parametrized by the second-moment tensor of a prescribed rectangular averaging window, and the memory is integrated by the L1 scheme in a reduced modal representation. Within this single framework we separate the effects of window geometry and orientation, electrode configuration, strip thickness, and the choice of stiffness operator that carries the memory.

The main results, all for the assumed BaTiO$_3$ parameters and the adopted model, are:
\begin{itemize}
\item the extensional coupling factor $k^{2}=0.5727$ agrees with the closed-form value $0.5752$ to $0.44\,\%$, and the quasi-static coupling is independent of thickness to better than $0.08\,\%$;
\item rotating an anisotropic window changes the extensional coupling by $14$--$17\,\%$;
\item reducing the thickness from $10$ to $1~\mu$m raises the flexoelectric bending coupling by a factor of about $24$ in the archived mesh study, while the piezoelectric extensional coupling is essentially unchanged;
\item the first-period equivalent damping depends on the Caputo order and, at fixed order, on whether the memory acts on the open-circuit or the mechanical-only stiffness (a ratio of $0.431$ at $\alpha=1$).
\end{itemize}

We state the scope of the work plainly in the manuscript. The gradient energy is an adopted local closure, the flexoelectric coefficient and the window parameters are prescribed rather than calibrated against experiment, part of the design map uses windows larger than the strip, and the very thin-strip results are not certified to high accuracy. The results therefore quantify the sensitivity of the adopted model and are not presented as calibrated device predictions.

The manuscript is original, has not been published previously and is not under consideration elsewhere. All authors have read and approved the submitted version and agree to its submission. The authors report no conflict of interest. The numerical results can be reproduced with the accompanying code and archived output files.

\textit{[Authors to confirm each statement in this paragraph and the funding acknowledgement before submission.]}

Thank you for considering our work. We look forward to hearing from you.

\medskip
Yours sincerely,

\medskip
Bandar Almohsen (corresponding author)\\
on behalf of Vipin Gupta, Sunita Kumawat, Arun Kumar Yadav and Taizo Maruyama\\
Department of Mathematics, College of Science, King Saud University, Riyadh, Saudi Arabia\\
balmohsen@ksu.edu.sa
\end{document}
